% META: {"id": "1SPE-SECDEG-EX-025", "chapitre": "1SPE-SECOND-DEGRE", "type_objet": "exercice", "capacites": ["1SPE-SECOND-DEGRE-C4", "1SPE-SECOND-DEGRE-C5"], "capacites_codes": ["C4", "C5"], "methodes": ["M4", "M5"], "parcours": 2, "competences": ["calculer", "raisonner"], "duree_min": 20, "mode_creation": "ex_nihilo", "sources_inspiration": [], "parametres_sympy": {}, "fichier_tex": "chapitres/1SPE-SECOND-DEGRE/exercices/1SPE-SECDEG-EX-025.tex", "corrige_tex": "chapitres/1SPE-SECOND-DEGRE/corriges/1SPE-SECDEG-CO-025.tex", "status": "generated"}
% BEGIN-VERIFY
% from sympy import symbols, solve, simplify, Rational
% x = symbols('x')
% # f(x) = x^2 - 3x - 10 = (x-5)(x+2), racines 5 et -2
% f = x**2 - 3*x - 10
% sols = sorted(solve(f, x))
% assert sols == [-2, 5]
% assert simplify(f - (x-5)*(x+2)) == 0
% # signe : a>0, negatif entre -2 et 5
% assert f.subs(x, 0) < 0
% assert f.subs(x, 6) > 0
% # inequation f(x) < 0 : solution ]-2 ; 5[
% # inequation f(x) >= 0 : solution ]-inf;-2] U [5;+inf[
% # g(x) = -x^2 + 2x + 8 = -(x^2-2x-8) = -(x-4)(x+2), racines 4 et -2
% g = -x**2 + 2*x + 8
% sols_g = sorted(solve(g, x))
% assert sols_g == [-2, 4]
% assert g.subs(x, 0) > 0   # entre -2 et 4
% assert g.subs(x, 5) < 0   # en dehors
% END-VERIFY
\begin{exercice}{1SPE-SECDEG-EX-025}{2}{20}
\begin{enumerate}
  \item \textbf{(C4)} Factoriser $f(x) = x^2 - 3x - 10$ et dresser son tableau de signes sur $\mathbb{R}$.
  \item \textbf{(C5)} Résoudre dans $\mathbb{R}$ :
  \begin{enumerate}
    \item $x^2 - 3x - 10 < 0$
    \item $x^2 - 3x - 10 \geqslant 0$
  \end{enumerate}
  \item \textbf{(C4)} Factoriser $g(x) = -x^2 + 2x + 8$ et dresser son tableau de signes.
  \item \textbf{(C5)} Résoudre dans $\mathbb{R}$ : $-x^2 + 2x + 8 > 0$.
\end{enumerate}
\end{exercice}
